\subsection*{How \texttt{.co} files compile}

A \texttt{.co} file mixes three kinds of code. The Croqtile compiler treats each differently --- and the difference matters for understanding what you can and cannot do at each boundary. Consider this skeleton:

\begin{lstlisting}[language=croqtile,caption={How .co files compile}]
__device__ __forceinline__ float fast_rsqrt(float x) {   // 1
  return __frsqrt_rn(x);
}
__co__ void my_kernel(f32 [M, N] input, f32 [M, N]& output) {   // 2
  parallel {bm, bn} by [cdiv(M, 64), cdiv(N, 64)] : block {
    // ... Croqtile orchestration: parallel, dma, mma, events ...
    __cpp__("asm volatile(\"setmaxnreg.dec.sync.aligned.u32 40;\");");   // 3
  }
}
int main() {   // 4
  auto in = croq::make_spandata<croq::f32>(M, N);
  auto out = my_kernel(in.view());
}
\end{lstlisting}

The compiler splits this into a single intermediate \texttt{.cu} file, then hands it to \texttt{c++}:

\begin{figure}[htbp]\centering
\includegraphics[width=0.86\linewidth,height=0.42\textheight,keepaspectratio]{f10-fig1_compilation_flow_light.png}
\end{figure}

What happens to each region:

\begin{itemize}[leftmargin=1.6em]
\item \textbf{(1) \texttt{\_\_device\_\_} functions} are copied into the generated \texttt{.cu} file \textbf{verbatim}. They appear in the intermediate source exactly as you wrote them.
\item \textbf{(2) \texttt{\_\_co\_\_} functions} are the heart of Croqtile.
\item The compiler does not parse their contents --- they must be valid at the splice point in the generated code.
\item \textbf{(4) Host code} (\texttt{main()} and any other non-annotated functions) is passed through as ordinary C++ and ends up in the host compilation path.
\end{itemize}

All four regions merge into one \texttt{.cu} file. Run \texttt{croqtile kernel.\allowbreak{}co -\allowbreak{}v} (Section~\ref{sec:1-5}) to see the exact \texttt{c++} command the compiler issues.
